% !TEX root = ../main.tex
\chapter{インターネットという分散網}
\label{chap:internet}

\begin{chapterabstract}
端末で組み立てた要求は、単一組織が管理する一本の通信路を進むわけではありません。
インターネットは、異なる管理主体のネットワークを相互接続した「網の網」です。
この章では、パケットの配送、組織間の経路交換、名前解決を分けて扱います。
\end{chapterabstract}

\begin{learninggoals}
  \item 通信性能を、帯域幅、遅延、往復時間、揺らぎ、損失に分けて説明できます。
  \item IP、ルーティング、AS、IX、BGPの役割を、同じものとして混同せず説明できます。
  \item DNSの委任とキャッシュを、\code{dig +trace}の結果へ対応づけられます。
\end{learninggoals}

\chapterimage[端末 → ルーター → 自律システム → 相互接続 → 複数経路 ／ DNS：名前を階層的に解決]{ch05-internet.png}{異なる管理主体の網が複数の経路で相互接続する概念図。右側の木は階層的な名前解決を表します。}{fig:ch05-internet}

\section{通信の要素}

離れた機器が規則に従ってデータを交換することを\term{通信}と呼びます。
データを出す側が\term{送信元}です。
データを受け取る側が\term{宛先}です。
信号を運ぶ物理的または論理的な経路が\term{通信路}です。
通信では、宛先の識別、途中での転送、到着後の受け渡しを分けて考えます。

\section{通信性能}

\term{遅延}は、送信を始めてから相手側で処理できるまでに要する時間です。
\term{帯域幅}は、通信路が単位時間に運べるデータ量の上限です。
\term{スループット}は、測定区間で実際に運べたデータ量です。
帯域幅が大きくても、待ち時間や損失があればスループットは下がります。

\term{RTT}はRound-Trip Timeの略で、送信元から宛先へ到達して返事が戻るまでの往復時間です。
\term{ジッター}は、複数回測った遅延の揺らぎです。
アプリケーションの体感は、帯域幅だけでは判断できません。

\section{回線交換とパケット交換}

\term{回線交換}は、通信前に経路上の容量を予約する方式です。
予約した容量は安定しやすい一方、送信していない時間も他の通信へ回しにくくなります。

\term{パケット交換}は、データを小さな単位へ分けて共有回線へ流す方式です。
\term{パケット}は、パケット交換でデータを分けて運ぶ小さな単位です。
各パケットは別の通信と同じ回線を時間的に共有できます。
インターネットは、パケット交換によって多様な通信を同じ基盤へ載せます。

\section{パケットの構造}

\term{ヘッダー}は、宛先や制御情報を置く先頭部分です。
\term{ペイロード}は、上位層から受け取った運搬対象です。
一つのリンクで運べる最大サイズを\term{MTU}と呼びます。
MTUに収まるよう分割された断片が\term{フラグメント}です。
IPv4では途中の分割があり得ますが、IPv6では送信元側のサイズ調整が前提です。

中継装置がすぐ送れないパケットを置く列が\term{待ち行列}です。
待ち行列で生じる時間が\term{キューイング遅延}です。
待ち行列が保持限界を超えたときにパケットが破棄される現象を\term{パケット損失}と呼びます。

\section{プロトコル}

\term{プロトコル}は、通信する主体が共有する形式、意味、進行手順の約束です。
形式だけ一致しても、応答する時点や失敗時の処理が違えば相互運用できません。
仕様を読むときは、メッセージの形、各フィールドの意味、状態遷移を分けて確認します。

\section{階層化とカプセル化}

\term{階層化}は、通信機能を役割ごとの層に分ける設計です。
\term{カプセル化}は、上位層から受け取ったデータをペイロードにし、自層のヘッダーを付ける処理です。
受信側は外側のヘッダーから順に読み、内側のデータを上位層へ渡します。

\section{OSIとTCP/IP}

\term{OSI参照モデル}は、通信機能を7層で整理する参照用のモデルです。
\term{TCP/IPモデル}は、インターネットで使われるプロトコル群を中心に層を整理します。
両者は厳密な一対一対応ではありません。
本書では、役割を説明する地図として使い、実装をモデルへ無理に当てはめません。

\begin{figure}[htbp]
  \centering
  \begin{tikzpicture}[x=1mm,y=1mm,every node/.append style={font=\scriptsize}]
    \node[sffamily\bfseries,text=DeepNavy] at (29,57) {OSI参照モデル};
    \node[sffamily\bfseries,text=DeepNavy] at (93,57) {TCP/IPモデル};
    \foreach \y/\n/\labeltext/\example in {
      48/7/アプリケーション層/HTTP・DNS,
      40/6/プレゼンテーション層/表現変換,
      32/5/セッション層/対話管理,
      24/4/トランスポート層/TCP・UDP,
      16/3/ネットワーク層/IP,
      8/2/データリンク層/Ethernet,
      0/1/物理層/電気・光・無線
    } {
      \node[box,minimum width=54mm,minimum height=7mm,anchor=south west] at (0,\y) {\n：\labeltext\quad{\tiny \example}};
    }
    \node[accentbox,minimum width=54mm,minimum height=23mm,anchor=south west] (app) at (66,32) {アプリケーション層\\{\tiny HTTP・DNSなど}};
    \node[accentbox,minimum width=54mm,minimum height=7mm,anchor=south west] (trans) at (66,24) {トランスポート層\quad{\tiny TCP・UDP}};
    \node[accentbox,minimum width=54mm,minimum height=7mm,anchor=south west] (internet) at (66,16) {インターネット層\quad{\tiny IP}};
    \node[accentbox,minimum width=54mm,minimum height=15mm,anchor=south west] (link) at (66,0) {リンク層\\{\tiny Ethernet・無線LANなど}};
    \draw[softflow] (54,51.5) -- (app.west);
    \draw[softflow] (54,43.5) -- (app.west);
    \draw[softflow] (54,35.5) -- (app.west);
    \draw[softflow] (54,27.5) -- (trans.west);
    \draw[softflow] (54,19.5) -- (internet.west);
    \draw[softflow] (54,11.5) -- (link.west);
    \draw[softflow] (54,3.5) -- (link.west);
  \end{tikzpicture}
  \caption{OSI参照モデルの7層とTCP/IPモデルの概念的な対応。対応は説明用であり、厳密な一対一ではありません。}
  \label{fig:osi-tcpip-models}
\end{figure}

OSIの第7層から第5層は、アプリケーション間の機能、データ表現、対話の管理を分けます。TCP/IPモデルでは、これらをアプリケーション層の内部で扱うことが一般的です。第4層は端点間の配送、第3層はネットワークをまたぐ配送、第2層と第1層は隣接機器間の転送と信号を扱います。

\source{\href{https://www.itu.int/rec/T-REC-X.200}{ITU-T X.200 / ISO/IEC 7498-1: OSI Basic Reference Model}、\href{https://www.rfc-editor.org/rfc/rfc1122.html}{RFC 1122: Requirements for Internet Hosts}}

\section{IPアドレスとCIDR}

\term{IP}はInternet Protocolの略で、ネットワーク間でデータグラムを配送する約束です。
\term{IPv4}は32ビットのアドレスを用いるIPです。
\term{IPv6}は128ビットのアドレスを用いるIPです。
\term{IPアドレス}は、IP層で送信元と宛先を識別する値です。

\term{サブネット}は、共通のネットワーク接頭辞を持つアドレス範囲です。
\term{CIDR}はClassless Inter-Domain Routingの略で、接頭辞長を用いて範囲を表します。
例えば\code{192.0.2.0/24}では、先頭24ビットがネットワーク部分です。

\section{ルーティング}

\term{ルーター}は、受け取ったパケットを次の通信路へ転送する装置または機能です。
\term{ルーティング}は、宛先に応じて次の転送先を選ぶ処理です。
\term{ルーティングテーブル}は、宛先範囲と次の転送先の対応を保持します。
\term{デフォルトゲートウェイ}は、より具体的な経路がない場合に端末が送る転送先です。

\term{NAT}はNetwork Address Translationの略で、アドレスやポートを変換する仕組みです。
NATはアドレス不足への対処や境界設計に使われますが、IPの暗号化や相手の認証を保証しません。

\section{網の網}

\term{ISP}はInternet Service Providerの略で、インターネット接続を提供する事業者です。
\term{自律システム}は、一つの運用方針で経路を管理するネットワーク群です。
自律システムはAutonomous Systemを略して\term{AS}と呼びます。

\term{IX}はInternet Exchangeの略です。
\term{インターネットエクスチェンジ}は、複数のネットワークが通信を交換するための接続基盤です。
インターネット全体を単一組織が管理するのではなく、AS同士が接続関係と方針を持ち寄ります。

\section{経路制御の自律性}

\term{BGP}はBorder Gateway Protocolの略で、AS間で到達可能な宛先範囲を交換します。
\term{経路広告}は、ある宛先へ到達できることを隣接ASへ知らせる情報です。
\term{AS\_PATH}は、経路広告が通過したASの並びです。
\term{到達性}は、ある宛先へパケットを届けられる性質です。

\term{経路制御}は、到達性だけでなく、契約や運用方針を踏まえて経路を選ぶ処理です。
したがって、AS\_PATHが最短の経路を必ず採用するわけではありません。
\source{\href{https://datatracker.ietf.org/doc/html/rfc4271}{RFC 4271: Border Gateway Protocol 4}}

\section{中央のない網の耐障害性}

\term{冗長性}は、同じ役割を果たせる経路や部品を複数持つ性質です。
\term{耐障害性}は、一部が故障しても必要な機能を継続できる性質です。
\term{単一障害点}は、その一箇所の故障だけで全体が停止する要素です。
\term{迂回}は、利用できない区間を避けて別の経路を選ぶことです。

IPは到着、順序、所要時間を保証しない\term{ベストエフォート}の配送です。
別経路があることは自動的な復旧を保証しません。
経路情報の伝播時間や誤った広告によって、一時的に到達できない場合があります。

\section{エンドツーエンド原則}

\term{エンドツーエンド原則}は、通信の正しさを最終的に判断できる端点へ機能を置く設計上の考え方です。
途中の網は汎用的なパケット配送へ集中します。
到着確認やアプリケーション固有の整合性は、必要に応じて端点側で実現します。
すべての機能を端点だけへ置くという絶対規則ではありません。

\section{DNS}

\term{DNS}はDomain Name Systemの略で、名前を階層的に管理する分散データベースです。
\term{ドメイン名}は、人やアプリケーションが参照する階層的な名前です。
\term{ルートサーバー}は、名前空間の最上位から次の問い合わせ先を案内します。
\term{TLD}はTop-Level Domainの略で、\code{com}や\code{jp}など最上位の区分です。

\term{権威DNS}は、担当範囲の正式なレコードを返すサーバーです。
\term{リゾルバ}は、利用者側に代わって名前解決を進めます。
\term{委任}は、下位の名前空間を別の権威DNSへ任せる仕組みです。
\term{キャッシュ}は、得た応答を再利用する保存です。
\term{TTL}はTime to Liveの略で、DNS応答を再利用できる時間の目安です。
\source{\href{https://datatracker.ietf.org/doc/html/rfc1034}{RFC 1034}、\href{https://datatracker.ietf.org/doc/html/rfc1035}{RFC 1035}}

\section{手を動かす：経路を観測する}

\term{traceroute}は、宛先までの中継点を応答可能な範囲で観測するコマンドです。
\term{dig}は、DNSへ問い合わせてレコードや委任を確認するコマンドです。
次の操作は、自組織の規則と対象サービスの利用条件を確認してから実行してください。

\begin{lstlisting}[language=bash]
traceroute example.com
dig +trace example.com
\end{lstlisting}

\code{traceroute}の各行は、必ず物理経路の全装置を表すわけではありません。
\code{dig +trace}では、ルートからTLD、権威DNSへ委任が移る順序を確認します。

\section{網から接続へ}

IP層で宛先まで運べても、同じ端末上では複数のアプリケーションが待っています。
次章では、宛先の端末内で通信を届け分け、順序や再送を加え、安全なHTTP要求へ組み上げます。

\section{旅をたどり直す}

注文確認AIの要求は、ドメイン名からIPアドレスを調べ、端末のルーティングテーブルに従って外へ出ます。
複数のASはBGPで交換した情報と運用方針を用いてパケットを転送します。
この段階で得られるのは、宛先付近まで運ぶベストエフォートの基盤です。
\takeaway{インターネットは、パケット配送、組織ごとの経路制御、階層的な名前解決を組み合わせた分散網です。}
